% chapters/07_future_directions.tex

\section{Overview}

This chapter discusses natural extensions of the framework developed in
Chapters 3--6. The extensions are grouped by mathematical character:
analytic extensions of the model (Section~\ref{sec:fut-analytic}),
numerical and computational extensions (Section~\ref{sec:fut-numerical}),
and formal-verification extensions (Section~\ref{sec:fut-lean}).

\section{Analytic Extensions}
\label{sec:fut-analytic}

\subsection{Jump-Diffusion Common Noise}

The current framework assumes continuous common noise modelled by a Brownian motion
$W^0$. A natural extension is to allow \emph{jump-diffusion common noise}:
\[
    dX_t = b(t, X_t, \mu_t, \alpha_t)\, dt + \sigma\, dW_t + \sigma_0\, dW_t^0
    + \int_{\mathbb{R}} \gamma(t, X_{t-}, \mu_{t-}, z)\, (\eta - \bar\eta)(dt, dz),
\]
where $\eta$ is a Poisson random measure with compensator $\bar\eta$. This models
sudden common shocks (e.g., macroeconomic announcements, regime changes, natural
disasters) affecting all agents simultaneously. The FBSDE characterisation
acquires a jump component in the adjoint equation, and the master equation gains
an integral term involving finite differences of $U$ in the measure argument.
Existence and uniqueness for this class was studied for the deterministic-$\mu$
case in \citep{carmona2013probabilistic}; the common-noise analogue remains open.
arXiv:2401.10952 (Bayraktar and Cohen, 2024) initiates this study in finite
state spaces; the continuous-state extension is a priority for future work.

\subsection{Mean-Field Control (Social Optimum)}

The current thesis studies \emph{mean-field games} (competitive equilibrium), where
each agent minimises its own cost. The complementary problem is \emph{mean-field
control} (social optimum), where a central planner minimises the aggregate cost
$\frac{1}{N}\sum_i J^i$. The social optimum satisfies a different FBSDE system in
which the interaction term carries an additional derivative with respect to the
measure. The gap between the competitive equilibrium and the social optimum --- the
\emph{price of anarchy} --- is an important quantity in multi-agent systems.
Quantifying this gap in the common-noise setting via a Wasserstein distance is
an open problem.

\subsection{Full Convergence via Lata\l{}a--Oleszkiewicz}

Theorem~\ref{thm:intro-wellposedness} establishes \emph{subsequential} convergence
of the iterates of the fixed-point map $\Phi$. Full convergence (not just along a
subsequence) follows from uniqueness of the equilibrium and the contraction theorem,
which are already established. However, the convergence of the reinforcement
learning actor-critic algorithm of earlier work is only subsequential, because
the limit-point theorem for stochastic approximation gives $\liminf W_2 = 0$
rather than $\lim W_2 = 0$. The Lata\l{}a--Oleszkiewicz inequality
\[
    W_2(\mu, \mu^*)^2 \le C \cdot \mathrm{Ent}(\mu \| \mu^*)
    \quad (\text{log-Sobolev class})
\]
combined with an entropy descent argument could convert the subsequential limit
to a full limit, at the cost of restricting to log-Sobolev policy classes. This
route remains open.

\subsection{McKean--Vlasov Games with Asymmetric Information}

All agents in the current model observe the same common noise $W^0$. Extending to
\emph{asymmetric information} --- where agent $i$ observes a private signal
$Y_t^i = h(X_t^i, W^0_t) + \varepsilon_t^i$ in addition to a common signal ---
leads to a filtering problem for each agent. The equilibrium measure $\mu_t$ then
becomes the conditional distribution given the agent's information set, and the
master equation acquires a Zakai-type SPDE term. This extension connects the MFG
framework to nonlinear filtering theory.

\section{Numerical and Computational Extensions}
\label{sec:fut-numerical}

\subsection{Deep Galerkin Methods}

The finite-difference scheme of Chapter 4 faces a \emph{curse of dimensionality}
for $d \ge 2$: the complexity scales exponentially with $d$. Deep Galerkin methods
(DGM, \citep{sirignano2018}) parametrise the value function $U(t,x,\mu)$ as a
neural network and minimise the residual of the master equation
\eqref{eq:master-equation}. The challenge is representing the measure argument
$\mu$ as an input; moment embeddings (using $(\bar\mu, v, \ldots)$) or the
Sinkhorn attention mechanism can serve as finite-dimensional surrogates. Error
analysis for DGM in the MFG setting is an active research area.

\subsection{Physics-Informed Neural Networks for the Fokker--Planck SPDE}

The Fokker--Planck SPDE
\[
    d\rho_t = -\nabla_x \cdot (\rho_t b_t)\, dt + \frac{\sigma^2}{2}\Delta_x \rho_t\, dt
    + \sigma_0 \nabla_x \cdot (\rho_t \circ dW_t^0)
\]
is a stochastic PDE in $(d+1)$ dimensions. Physics-informed neural networks
(PINNs) can enforce the PDE residual in the loss function. For the common-noise
case, the stochastic term requires an It\^{o}--Stratonovich correction; implementing
this in an automatic-differentiation framework (JAX, PyTorch) is straightforward
but has not been systematically studied.

\subsection{High-Dimensional Common-Noise MFGs}

The particle method of Chapter 5 achieves $O(N^{-2/d})$ rate for $d \ge 5$,
degrading rapidly with dimension. Three approaches are known to mitigate this:
(a) \emph{structured state spaces} where the interaction is through low-dimensional
statistics (e.g., first two moments); (b) \emph{normalising flows} that map the
measure to a parametric family; (c) \emph{diffusion models} (score-based
generative models) that directly approximate $\mu_t^*$ from samples. None of these
has a rigorous convergence theory in the common-noise MFG setting.

\subsection{Matrix Riccati Re-derivation for Chapter 8's Multi-Dimensional LQ-MFG}

The corrected Appendix A derivation covers the scalar LQ-MFG (Chapter 7 of the
original thesis). Chapter 8's alpha-generation model uses a matrix-valued Riccati
system, where the same sign corrections likely apply but have not been independently
re-derived. Providing a verified matrix re-derivation analogous to Appendix A ---
including a symbolic PDE-residual check and numerical solver cross-validation ---
is the primary technical debt left by the revision process and is identified as the
first priority for future work.

\section{Formal Verification Extensions}
\label{sec:fut-lean}

\subsection{Full FBSDE Well-Posedness in Lean 4}

The current Lean 4 proofs (Chapter 6) cover the key analytic lemmas but not the
full FBSDE well-posedness theorem. A complete formalisation would require:
\begin{enumerate}
    \item Measure-valued SDEs in Lean / Mathlib (not yet available as of Mathlib
    4.14; relevant PRs are in progress);
    \item The Banach fixed-point theorem applied to the function space
    $\mathcal{C}([0,T]; \mathcal{P}_2(\mathbb{R}^d))$ (available via
    \texttt{Mathlib.Topology.MetricSpace.Contracting});
    \item Measurability of the conditional measure flow $t \mapsto \mu_t$ with
    respect to $\mathcal{F}_t^{W^0}$ (requires stochastic analysis infrastructure
    in Lean not yet in Mathlib).
\end{enumerate}

\subsection{Verified Numerical Scheme}

A longer-term goal is a Lean 4 / Coq proof of the convergence theorem
(Theorem~\ref{thm:intro-numerical}) for the Euler--Maruyama scheme. This requires
a verified floating-point arithmetic model and interval arithmetic to bound
rounding errors. The Lean 4 formalisation of Euler's method for ODEs in
\texttt{Mathlib.Analysis.ODE.Picard} provides a starting point, but
extension to stochastic integrals and Wasserstein distances requires substantial
new infrastructure.

\section{Conclusion}

The framework developed in this thesis --- quantitative well-posedness
(Chapter 3), convergent numerical schemes (Chapter 4), propagation-of-chaos
analysis (Chapter 5), and partial formal verification (Chapter 6) --- provides
a rigorous foundation for the mathematical study of interacting particle systems
under common noise. The extensions described in this chapter map out a research
agenda that spans stochastic analysis, computational mathematics, and formal
methods: exactly the intersection where the most interesting open problems
currently lie.
